% Ad Atticum 7.1 --- headnote
% ad-atticum-07-01, 16 October 50 BC, Athens

\textit{Cicero to Atticus, written from Athens on the
seventeenth day before the Kalends of November 50 BC ---
16 October --- the opening letter of book 7 of the
\textit{Ad Atticum}, and the last surviving letter from
the eastward voyage home from Cilicia. He had landed at
the Piraeus on the 14th, written briefly the next day
(book 6.9) and sent that off with Saufeius, and now ---
two days later, before sailing for Italy --- sits down
to the long, anxious letter he could not fit in the
first time. The seal between books 6 and 7 is editorial,
not biographical: the same voice continues.}

\textit{Section 1 is the recap, deliberately laid down
for the case in which the Saufeius letter does not arrive
first --- philosophers, as he says, do not walk fast.
Sections 2--5 are the heart: the prayer for Atticus's
\textit{prudentia}, the open citation of his own
miscalculation in joining both Caesar and Pompey
(\textit{ut neutri illorum quisquam esset me carior},
``so that neither of them held anyone dearer than me''),
and the rehearsal --- in mock-Senate procedural form ---
of the impossible vote he sees coming: \textit{dic, M.
Tulli. quid dicam?} ``Speak, Marcus Tullius. What shall
I say?'' The Greek thickens accordingly: a Homeric
half-line about a heart never persuaded
(\emph{ἀλλ' ἐμὸν οὔποτε θυμὸν ἐνὶ στήθεσσιν ἔπειθεσ},
\textit{Iliad} 9.345 of Achilles to the embassy), an
\emph{αἰδέομαι} drawn from Hector's speech in
\textit{Iliad} 22, and \emph{Πουλυδάμας μοι} ---
``Polydamas first of all'' --- the first words of Hector's
self-reproach for refusing prudent counsel. He breaks off
to answer his own question: who is the Polydamas?
\textit{tu ipse} --- ``you yourself.''}

\textit{Sections 6--9 turn back to the manageable
business: the provincial accounts (his unfashionable
honesty in returning the surplus to the treasury, against
the muttering of his own \textit{cohors}), the
thanksgiving and the triumph he hopes for, the lobbying of
Hirrus, and the domestic Precius affair already
glimpsed in 6.9 --- the same \emph{φυρατής}, ``schemer,''
now named outright as ``Lartidius to the life.'' The
letter ends as it began: by handing Atticus the work to
do. The Rubicon is two and a half months away.}
